% !TeX program = lualatex
% ================================================================
% Polish Tier 3 Vocabulary Version of SequencesStartersPatterns
%
% Generated by the server using the following settings:
%
%   language            Polish (pol)
%   text direction      left to right
%   font                Noto Serif
%   fallback font       Noto Serif
%   built from          latex/starters/targeted/KS3/sequences/sequencesStartersPatterns.tex
%   vocabulary key      latex/starters/targeted/KS3/sequences/sequencesStartersPatterns/L2/pol/sequencesStartersPatterns_polishKey.tex
%   tier 3 vocabulary   green   RGB 0 128 0
%   English text        black   RGB 0 0 0
%   Polish text         purple  RGB 112 48 160
%   layout macros       version 3
%
% Please don't edit any information in this comment block - the server
% compares it against the current settings and rebuilds this file when
% they differ, so changes here would be overwritten. However, feel free
% to make updates to the LaTeX below if you want to edit it.
% ================================================================

\documentclass{beamer}
% ================================================================
% Copied from the English sheet this was translated from, so that
% anything it sets up for itself still works here
% ================================================================
\usepackage{tikz}
\usepackage{multicol}

% ================================================================
% Answer-overlay helpers
% ================================================================
\newcommand{\ablank}[1]{%
  \alt<2>{\textcolor{red}{#1}}{\underline{\phantom{#1}}}%
}
\newcommand{\ashow}[1]{\uncover<2->{\textcolor{red}{\small #1}}}
\newcommand{\ashowq}[1]{\alt<2->{\textcolor{red}{\small #1}}{?}}

% Answers drawn straight on to a TikZ diagram, revealed on overlay 2 like \ashow.
% Space is reserved on overlay 1, so the diagram does not move between slides.
\usetikzlibrary{overlay-beamer-styles}
\tikzset{
  ans/.style     = {red, thick, visible on=<2->},
  ansfill/.style = {red!15, visible on=<2->},
  anslab/.style  = {red, font=\tiny, visible on=<2->},
}

% ---- Minimal styling ----
\setbeamertemplate{navigation symbols}{}
\setbeamertemplate{footline}{}
\setbeamercolor{background canvas}{bg=white}
\setbeamercolor{frametitle}{fg=black}
\setbeamerfont{frametitle}{size=\Large, series=\bfseries}
\usefonttheme{professionalfonts}

% ---- Shared tikz ----
\tikzset{every picture/.style={baseline=(current bounding box.south)}}

% 1) matchstick squares (n squares)
\newcommand{\matchsquares}[2][black]{%
  \begin{tikzpicture}[scale=0.65]
    \pgfmathsetmacro\n{#2}
    \color{#1}
    \foreach \x in {0,...,\n} {
      \draw[very thick, cap=round] (\x,0) -- (\x,1);
    }
    \foreach \x in {0,...,\the\numexpr\n-1\relax} {
      \draw[very thick, cap=round] (\x,0) -- (\x+1,0);
      \draw[very thick, cap=round] (\x,1) -- (\x+1,1);
    }
  \end{tikzpicture}%
}

% 2) dot squares (n x n)
\newcommand{\dotgrid}[2][black]{%
  \begin{tikzpicture}[scale=0.65]
    \pgfmathsetmacro\n{#2}
    \color{#1}
    \foreach \i in {0,...,\the\numexpr\n-1\relax} {
      \foreach \j in {0,...,\the\numexpr\n-1\relax} {
        \fill (0.4*\i,0.4*\j) circle (0.07);
      }
    }
  \end{tikzpicture}%
}

% 3) matchstick triangles (chain)
\newcommand{\matchtriangles}[2][black]{%
  \begin{tikzpicture}[scale=0.65]
    \pgfmathsetmacro\n{#2}
    \pgfmathsetmacro\h{0.866}
    \color{#1}
    \draw[very thick] (0,0) -- (\n,0);
    \draw[very thick] (0,0)
      \foreach \k in {1,...,\n} {
        -- ({\k - 0.5}, {\h}) -- (\k,0)
      };
  \end{tikzpicture}%
}

% 4) triangular numbers (dots)
\newcommand{\tridots}[2][black]{%
  \begin{tikzpicture}[scale=0.7]
    \pgfmathsetmacro\n{#2}
    \color{#1}
    \foreach \row in {1,...,\n} {
      \pgfmathsetmacro\k{int(\n - \row + 1)}
      \pgfmathsetmacro\xstart{ -0.45*(\k-1)/2 }
      \foreach \col in {0,...,\the\numexpr\k-1\relax} {
        \pgfmathsetmacro\x{\xstart + 0.45*\col}
        \pgfmathsetmacro\y{0.4*(\row-1)}
        \fill (\x,\y) circle (0.07);
      }
    }
  \end{tikzpicture}%
}

% 5) rectangle dots (n x (n+1))
\newcommand{\rectdots}[2][black]{%
  \begin{tikzpicture}[scale=0.65]
    \pgfmathsetmacro\n{#2}
    \color{#1}
    \foreach \i in {0,...,\the\numexpr\n-1\relax} {
      \foreach \j in {0,...,\n} {
        \fill (0.4*\i, 0.4*\j) circle (0.07);
      }
    }
  \end{tikzpicture}%
}

% 6) cross dots
\newcommand{\crossdots}[2][black]{%
  \begin{tikzpicture}[scale=0.65]
    \pgfmathsetmacro\n{#2}
    \color{#1}
    \fill (0,0) circle (0.07);
    \ifnum\n>1
      \foreach \d in {1,...,\the\numexpr\n-1\relax} {
        \fill (0, 0.4*\d) circle (0.07);
        \fill (0, -0.4*\d) circle (0.07);
        \fill (0.4*\d, 0) circle (0.07);
        \fill (-0.4*\d, 0) circle (0.07);
      }
    \fi
  \end{tikzpicture}%
}

% Characters this language's font has not got are borrowed from another,
% rather than being silently dropped - see docs/translated-sheet-latex.md
\directlua{%
  if luaotfload and luaotfload.add_fallback then
    luaotfload.add_fallback("eallatinfallback",
      { "Noto Serif:mode=harf;" })
  else
    tex.error("this luaotfload is too old to register a fallback font")
  end
}
\usepackage[english]{babel}
\babelprovide[import]{polish}
\babelfont[polish]{rm}[Renderer=Harfbuzz, BoldFont={Noto Serif}, ItalicFont={Noto Serif}, BoldItalicFont={Noto Serif}, RawFeature={fallback=eallatinfallback}]{Noto Serif}
\babelfont[polish]{sf}[Renderer=Harfbuzz, BoldFont={Noto Serif}, ItalicFont={Noto Serif}, BoldItalicFont={Noto Serif}, RawFeature={fallback=eallatinfallback}]{Noto Serif}
\newcommand{\ealtext}[1]{\foreignlanguage{polish}{#1}}
\newcommand{\ealtextblock}[1]{{\raggedright\foreignlanguage{polish}{#1}\par}}
% ================================================================
% EAL layout helpers
% ================================================================
\usepackage{xcolor}

\definecolor{ealkeycolour}{RGB}{0,128,0}
\definecolor{ealenglish}{RGB}{0,0,0}
\definecolor{ealtwocolour}{RGB}{112,48,160}

\newsavebox{\ealboxen}
\newsavebox{\ealboxtr}
\newlength{\ealglosswd}
\newlength{\ealglossraise}
\setlength{\ealglossraise}{1.15em}

% a tier 3 word, in English and in translation alike
\newcommand{\ealkey}[1]{\textcolor{ealkeycolour}{#1}}
\newcommand{\ealkeytr}[1]{\textcolor{ealkeycolour}{\ealtext{#1}}}

% One translated word above one English word. Built from kernel
% primitives, never a tabular, which would break wherever the sheet
% loads array, booktabs or tabularx. The box is as wide as the wider of
% the two, so a long translation pushes its neighbours aside instead of
% overlapping them, and it claims its full height so a table row or a
% TikZ node makes room for it.
\newcommand{\ealgloss}[2]{%
  \sbox{\ealboxen}{\ealkey{#1}}%
  \sbox{\ealboxtr}{\small\textcolor{ealtwocolour}{\ealtext{#2}}}%
  \setlength{\ealglosswd}{\wd\ealboxen}%
  \ifdim\wd\ealboxtr>\ealglosswd\setlength{\ealglosswd}{\wd\ealboxtr}\fi
  \leavevmode
  \makebox[\ealglosswd][c]{%
    \makebox[0pt][c]{\usebox{\ealboxen}}%
    \makebox[0pt][c]{\raisebox{\ealglossraise}{\usebox{\ealboxtr}}}%
  }%
}

% opens up line spacing around a block of glosses, so the raised
% translations sit clear of the line above
\newenvironment{ealglossed}{\par\linespread{2.1}\selectfont\ignorespaces}{\par}

% a whole translated sentence above its English counterpart. block level,
% so it does not belong inside a tabular cell
\newcommand{\ealpara}[2]{%
  \par\addvspace{0.5\baselineskip}%
  {\color{ealtwocolour}\ealtextblock{#1}}%
  \penalty200
  {\color{ealenglish}#2\par}%
  \addvspace{0.5\baselineskip}%
}

\begin{document}

% =============================================
% TITLE
% =============================================
\begin{frame}[plain]
  \begin{ealglossed}
  \centering
  {\LARGE\bfseries \ealgloss{Sequences}{ciągi} Starters}\\[0.5em]
  \end{ealglossed}
\end{frame}

% =============================================
% STARTER 1 – Easy
% =============================================
\begin{frame}{Starter 1}
  \begin{minipage}[t]{0.60\textwidth}
    \centering
    \textbf{Matchsticks}\\[0.7em]
    \matchsquares[red]{1}\;\hspace{0.5cm}
    \matchsquares[blue]{2}\;\hspace{0.5cm}
    \matchsquares[green!70!black]{3}\;\hspace{0.5cm}
    \textbf{?}\\[0.8em]
    \small
    1. Draw the next shape.\\[0.2em]
    2. How many matchsticks in the 4th?
  \end{minipage}\hfill
  \begin{minipage}[t]{0.34\textwidth}
    \centering
    \textbf{Dots}\\[0.7em]
    \dotgrid[red]{1}\;\hspace{0.5cm}
    \dotgrid[blue]{2}\;\hspace{0.5cm}
    \dotgrid[green!70!black]{3}\;\hspace{0.5cm}
    \textbf{?}\\[0.8em]
    \small
    3. Draw the next pattern.\\[0.2em]
    4. How many dots in the 4th?
  \end{minipage}

  \vspace{0.8cm}                     % increased vertical space
  \begin{ealglossed}
  \centering
  \textbf{5. Complete The \ealgloss{Sequence}{ciąg}}\\[0.3em]
  \small
  13,\;23,\;33,\;
  \ablank{43},\;
  \ablank{53},\;
  \ablank{63}
  \ashow{\ (add 10)}
  \end{ealglossed}

  \vspace{1.2cm}                     % space before pattern answers
  \pause
  \begin{minipage}[t]{0.60\textwidth}
    \centering
    \textcolor{red}{\textbf{Answers}}\\[0.2em]
    \textcolor{red}{1. \matchsquares[orange]{4}}\\[0.2em]
    \textcolor{red}{2. 13 sticks}
  \end{minipage}\hfill
  \begin{minipage}[t]{0.34\textwidth}
    \centering
    \textcolor{red}{\textbf{Answers}}\\[0.2em]
    \textcolor{red}{3. \dotgrid[orange]{4}}\\[0.2em]
    \textcolor{red}{4. 16 dots}
  \end{minipage}
\end{frame}

% =============================================
% STARTER 2 – Linear patterns
% =============================================
\begin{frame}{Starter 2}
  \begin{minipage}[t]{0.60\textwidth}
    \centering
    \textbf{Matchsticks}\\[0.7em]
    \matchtriangles[red]{1}\;\hspace{0.5cm}
    \matchtriangles[blue]{2}\;\hspace{0.5cm}
    \matchtriangles[green!70!black]{3}\;\hspace{0.5cm}
    \textbf{?}\\[0.8em]
    \small
    1. Draw shape 4.\\[0.2em]
    2. How many matchsticks in shape 4?
  \end{minipage}\hfill
  \begin{minipage}[t]{0.34\textwidth}
    \centering
    \textbf{Dots}\\[0.7em]
    \tridots[red]{1}\;\hspace{0.5cm}
    \tridots[blue]{2}\;\hspace{0.5cm}
    \tridots[green!70!black]{3}\;\hspace{0.5cm}
    \textbf{?}\\[0.8em]
    \small
    3. Draw pattern 4.\\[0.2em]
    4. How many dots in pattern 4?
  \end{minipage}

  \vspace{0.8cm}
  \begin{ealglossed}
  \centering
  \textbf{5. Complete The \ealgloss{Sequence}{ciąg}}\\[0.3em]
  \small
  1,\;3,\;5,\;7,\;
  \ablank{9},\;
  \ablank{11},\;
  \ablank{13}
  \ashow{\ (add 2)}
  \end{ealglossed}

  \vspace{1.2cm}
  \pause
  \begin{minipage}[t]{0.60\textwidth}
    \centering
    \textcolor{red}{\textbf{Answers}}\\[0.2em]
    \textcolor{red}{1. \matchtriangles[orange]{4}}\\[0.2em]
    \textcolor{red}{2. 9 sticks}
  \end{minipage}\hfill
  \begin{minipage}[t]{0.34\textwidth}
    \centering
    \textcolor{red}{\textbf{Answers}}\\[0.2em]
    \textcolor{red}{3. \tridots[orange]{4}}\\[0.2em]
    \textcolor{red}{4. 10 dots}
  \end{minipage}
\end{frame}

% =============================================
% STARTER 3 – Predicting further terms
% =============================================
\begin{frame}{Starter 3}
  \begin{minipage}[t]{0.36\textwidth}
    \centering
    \textbf{Dots}\\[0.6em]
    \rectdots[red]{1}\;\hspace{0.5cm}
    \rectdots[blue]{2}\;\hspace{0.5cm}
    \rectdots[green!70!black]{3}\;\hspace{0.5cm}
    \textbf{?}\\[0.8em]
    \small
    1. Draw pattern 4.\\[0.2em]
    2. How many dots in pattern 4?
  \end{minipage}\hfill
  \begin{minipage}[t]{0.60\textwidth}
    \begin{ealglossed}
    \centering
    \textbf{Matchsticks}\\[0.7em]
    \matchsquares[red]{1}\;\hspace{0.5cm}
    \matchsquares[blue]{2}\;\hspace{0.5cm}
    \matchsquares[green!70!black]{3}\;\hspace{0.5cm}
    \textbf{?}\\[0.8em]
    \small
    3. How many sticks in the 10th shape?\\[0.2em]
    4. Find a \ealgloss{rule}{reguła}.
    \end{ealglossed}
  \end{minipage}

  \vspace{0.8cm}
  \begin{ealglossed}
  \centering
  \textbf{5. Complete The \ealgloss{Sequence}{ciąg}}\\[0.3em]
  \small
  3,\;6,\;9,\;12,\;
  \ablank{15},\;
  \ablank{18},\;
  \ablank{21}
  \ashow{\ (add 3)}
  \hspace{1em} \footnotesize 6. 20th \ealgloss{term}{wyraz}? \ashow{$60$}
  \end{ealglossed}

  \vspace{0.5cm}
  \pause
  \begin{minipage}[t]{0.36\textwidth}
    \centering
    \textcolor{red}{\textbf{Answers}}\\[0.2em]
    \textcolor{red}{1. \rectdots[orange]{4}}\\[0.2em]
    \textcolor{red}{2. 20 dots \footnotesize (pattern is $n\times(n+1)$ dots)}
  \end{minipage}\hfill
  \begin{minipage}[t]{0.60\textwidth}
    \begin{ealglossed}
    \centering
    \textcolor{red}{\textbf{Answers}}\\[0.2em]
    \textcolor{red}{3. 31 }\\[0.2em]
    \textcolor{red}{4. \ealgloss{Rule}{reguła}: $3n+1$}
    \end{ealglossed}
  \end{minipage}
\end{frame}

% =============================================
% STARTER 4 – Nth term rules
% =============================================
\begin{frame}{Starter 4}
  \begin{minipage}[t]{0.60\textwidth}
    \begin{ealglossed}
    \centering
    \textbf{Matchsticks}\\[0.7em]
    \matchtriangles[red]{1}\;\hspace{0.5cm}
    \matchtriangles[blue]{2}\;\hspace{0.5cm}
    \matchtriangles[green!70!black]{3}\;\hspace{0.5cm}
    \textbf{?}\\[0.8em]
    \small
    1. Find a \ealgloss{rule}{reguła} for the number of matchsticks.\\[0.2em]
    2. Use it to find the 20th \ealgloss{term}{wyraz}.
    \end{ealglossed}
  \end{minipage}\hfill
  \begin{minipage}[t]{0.34\textwidth}
    \begin{ealglossed}
    \centering
    \textbf{Dots}\\[0.7em]
    \tridots[red]{1}\;\hspace{0.5cm}
    \tridots[blue]{2}\;\hspace{0.5cm}
    \tridots[green!70!black]{3}\;\hspace{0.5cm}
    \textbf{?}\\[0.8em]
    \small
    3. Find a \ealgloss{rule}{reguła} for the number of dots.\\[0.2em]
    4. Use it to find the 10th \ealgloss{term}{wyraz}.
    \end{ealglossed}
  \end{minipage}

  \vspace{0.8cm}
  \begin{ealglossed}
  \centering
  \textbf{5. Complete The \ealgloss{Sequence}{ciąg}}\\[0.3em]
  \small
  5,\;9,\;13,\;17,\;
  \ablank{21},\;
  \ablank{25},\;
  \ablank{29}
  \ashow{\ (add 4)}
  \hspace{1em} \footnotesize 6. \ealgloss{$n$th term}{wzór na n-ty wyraz} \ealgloss{rule}{reguła}? 100th \ealgloss{term}{wyraz}? \\ \hspace{12em} \ashow{$4n+1$, \quad \qquad \quad $401$}
  \end{ealglossed}

  \vspace{1.2cm}
  \pause
  \begin{minipage}[t]{0.60\textwidth}
    \begin{ealglossed}
    \centering
    \textcolor{red}{\textbf{Answers}}\\[0.2em]
    \textcolor{red}{1. \ealgloss{Rule}{reguła}: $2n+1$}\\[0.2em]
    \textcolor{red}{2. 20th \ealgloss{term}{wyraz}: $41$}
    \end{ealglossed}
  \end{minipage}\hfill
  \begin{minipage}[t]{0.34\textwidth}
    \begin{ealglossed}
    \centering
    \textcolor{red}{\textbf{Answers}}\\[0.2em]
    \textcolor{red}{3. \ealgloss{Rule}{reguła}: $\frac{n(n+1)}{2}$}\\[0.2em]
    \textcolor{red}{4. 10th \ealgloss{term}{wyraz}: $55$}
    \end{ealglossed}
  \end{minipage}
\end{frame}

% =============================================
% STARTER 5 – Challenge
% =============================================
\begin{frame}{Starter 5}
  \begin{minipage}[t]{0.50\textwidth}
    \begin{ealglossed}
    \centering
    \textbf{Dots}\\[0.7em]
    \crossdots[red]{1}\;\hspace{0.5cm}
    \crossdots[blue]{2}\;\hspace{0.5cm}
    \crossdots[green!70!black]{3}\;\hspace{0.5cm}
    \textbf{?}\\[0.8em]
    \small
    1. Draw shape 4.\\[0.2em]
    2. How many dots in shape 5?\\[0.2em]
    3. Try to find a \ealgloss{rule}{reguła}.
    \end{ealglossed}
  \end{minipage}\hfill
  \begin{minipage}[t]{0.50\textwidth}
    \begin{ealglossed}
    \centering
    \textbf{Special \ealgloss{Sequence}{ciąg}}\\[0.7em]
    1,\;1,\;2,\;3,\;5,\;8,\;
    \ablank{13},\;
    \ablank{21},\;
    \ablank{34}\\[0.8em]
    \small
    4. Next three \ealgloss{terms}{wyrazy}?\\[0.2em]
    5. What is the \ealgloss{rule}{reguła}?
    \end{ealglossed}
  \end{minipage}

  \vspace{0.8cm}
  \begin{ealglossed}
  \centering
  \textbf{6. Complete The \ealgloss{Sequence}{ciąg}}\\[0.3em]
  \small
  1,\;4,\;9,\;16,\;
  \ablank{25},\;
  \ablank{36},\;
  \ablank{49}
  \ashow{\ (\ealgloss{square numbers}{liczby kwadratowe} $n^2$)}
  \hspace{1em} \footnotesize 7. 20th \ealgloss{term}{wyraz}? \ashow{$400$}
  \end{ealglossed}

  \vspace{0.5cm}
  \pause
  \begin{minipage}[t]{0.50\textwidth}
    \begin{ealglossed}
    \begin{multicols}{2}
    \textcolor{red}{\textbf{Answers}}\\[0.2em]
    \textcolor{red}{1. \crossdots[orange]{4}}\\[0.2em]
    \textcolor{red}{2. Shape 5 has 17 dots}\\[0.2em]
    \textcolor{red}{3. \ealgloss{Rule}{reguła}: $4n-3$}
    \end{multicols}
    \end{ealglossed}
  \end{minipage}\hfill
  \begin{minipage}[t]{0.50\textwidth}
    \begin{ealglossed}
    \centering
    \vspace{-6.5em}
    \textcolor{red}{\textbf{Answers}}\\[0.2em]
    \textcolor{red}{4. 13,\;21,\;34}\\[0.2em]
    \textcolor{red}{5. \ealgloss{Fibonacci}{ciąg Fibonacciego}: add the two previous \ealgloss{terms}{wyrazy}}
    \end{ealglossed}
  \end{minipage}
\end{frame}

\end{document}